% basic-definitions.tex — \input'd from proofs.tex
% Conventions: thesis/CLAUDE.md
%
% Sources: Jörn's dictation + HK2017 paper. Notation per correspondence.tex.
%
% Structure: Fundamental definitions (symplectic form through shoelace formula)

% --- Smooth setting ---

The following section reintroduces definitions notation from symplectic geometry that the reader is assumed to be familiar with, but that may differ between sources.
Afterwards, we extend the usually smooth setting of symplectic geometry to polytopes, in particular we introduce the notion of generalized Reeb orbits.
When generalizing, we introduce as definitions the natural limit objects when going from smooth convex bodies to non-smooth polytopes.
See \cite{CH2021} for more detailed discussion and proofs that e.g. generalized Reeb orbits are indeed the limit of smooth Reeb orbits on smoothings of the polytope.

\begin{definition}[Standard symplectic structure]\label{def:symplectic-form}
% Jörn: text approved (1aac085)
% QC: J_0^2 = -I. omega_0 antisymmetric, nondegenerate.
%   omega_0(e_{q_i}, e_{p_i}) = +1 for i=1,2; all other basis pairs give 0.
% Downstream: J_0 is the 4x4 matrix [[0,-I_2],[I_2,0]].
%   omega_0(u,v) = dot(J_0*u, v) = u_q1*v_p1 - u_p1*v_q1 + u_q2*v_p2 - u_p2*v_q2.
%   Properties: J_0^2 = -I, J_0^T = J_0^{-1} = -J_0,
%   omega_0(J_0 u, J_0 v) = omega_0(u,v), |J_0 u| = |u|.
Write $\mathbb{R}^4$ with coordinates $(q_1, q_2, p_1, p_2)$.
The \emph{standard complex structure} is
\[
  J_0 = \begin{pmatrix} 0 & -I_2 \\ I_2 & 0 \end{pmatrix},
  \qquad\text{i.e.,}\quad
  J_0(q, p) = (-p, q),
\]
where $I_2$ is the $2 \times 2$ identity matrix.
The \emph{standard symplectic form} is
\[
  \omega_0(u, v) \coloneqq \langle J_0 u, v \rangle
  = \sum_{i=1}^{2} (u_{q_i} v_{p_i} - u_{p_i} v_{q_i}),
\]
equivalently $\omega_0 = dq_1 \wedge dp_1 + dq_2 \wedge dp_2$.
\end{definition}

\begin{remark}[Properties of $J_0$ and $\omega_0$]\label{rem:j0-properties}
% Jörn: text approved (1aac085)
% Downstream: all testable via proptest on random vectors.
~
\begin{enumerate}
\item $J_0^2 = -I_4$, \quad $J_0^{-1} = J_0^\top = -J_0$.
\item $\omega_0(J_0 u, J_0 v) = \omega_0(u, v)$ for all $u, v \in \mathbb{R}^4$.
\item $|J_0 u| = |u|$ for all $u \in \mathbb{R}^4$.
\item $\omega_0$ is nondegenerate: $\omega_0(u, v) = 0$ for all $v$ implies $u = 0$.
\end{enumerate}
\end{remark}

\begin{definition}[Liouville form]\label{def:liouville-form}
% Jörn: text approved (1aac085)
% QC: d(lambda_0) = omega_0.
%   lambda_0 at x applied to v equals (1/2) omega_0(x, v) = (1/2) <J_0 x, v>.
% Downstream: lambda_0(v) at point x = 0.5 * dot(J_0*x, v).
%   For a closed curve gamma, A(gamma) = integral of lambda_0 along gamma.
The \emph{standard Liouville form} is
\[
  \lambda_0 \coloneqq \tfrac{1}{2} \sum_{i=1}^{2} (q_i \, dp_i - p_i \, dq_i).
\]
At a point $x = (q, p) \in \mathbb{R}^4$, applied to a tangent vector $v$:
\[
  \lambda_0|_x(v) = \tfrac{1}{2}\, \omega_0(x, v) = \tfrac{1}{2}\, \langle J_0 x, v \rangle.
\]
It satisfies $d\lambda_0 = \omega_0$.
\end{definition}


\begin{definition}[Convex body]\label{def:convex-body}
% Jörn: text approved (84db708)
% QC: K compact, convex, nonempty interior, 0 in int(K).
% Downstream: the input to all algorithms. Always centered at origin.
A \emph{convex body} is a compact convex subset $K \subset \mathbb{R}^4$
with $0 \in \operatorname{int}(K)$.
\end{definition}

\begin{definition}[Support function]\label{def:support-function}
% Jörn: text approved (1aac085)
% QC: h_K positively 1-homogeneous, convex, h_K(v) > 0 for v != 0 (since 0 in int(K)).
% Downstream: h_K(v) = max_{x in K} dot(x, v).
%   For polytopes: h_K(n_i) = h_i. Test: h_K(alpha*v) = alpha * h_K(v) for alpha > 0.
The \emph{support function} of a convex body $K$ is $h_K \colon \mathbb{R}^4 \to \mathbb{R}$,
\[
  h_K(v) \coloneqq \max_{x \in K} \langle x, v \rangle.
\]
Since $0 \in \operatorname{int}(K)$, we have $h_K(v) > 0$ for all $v \neq 0$.
\end{definition}

\begin{definition}[Gauge function]\label{def:gauge-function}
% Jörn: text approved (1aac085)
% QC: g_K positively 1-homogeneous, convex, {g_K = 1} = bdry(K),
%   g_K < 1 on int(K), g_K > 1 outside K.
% Downstream: g_K(x) = inf { t > 0 : x in t*K }.
%   For polytopes: g_K(x) = max_i <x, n_i>/h_i.
%   Test: g_K(alpha*x) = alpha * g_K(x) for alpha > 0; g_K(x) = 1 iff x in bdry(K).
The \emph{gauge function} (Minkowski functional) of a convex body $K$ with
$0 \in \operatorname{int}(K)$ is $g_K \colon \mathbb{R}^4 \to [0,\infty)$,
\[
  g_K(x) \coloneqq \inf\{ \lambda > 0 : x \in \lambda K \}.
\]
The boundary is the level set: $\partial K = \{x : g_K(x) = 1\}$.
\end{definition}

\begin{definition}[Hamiltonian]\label{def:hamiltonian}
% Jörn: text approved (84db708)
% QC: H = g_K^2 per correspondence.tex. {H = 1} = bdry(K).
% Downstream: H(x) = g_K(x)^2. On interior of facet F_i of a polytope:
%   grad(H)(x) = 2 n_i / h_i, so J_0 grad(H) = R_i = (2/h_i) J_0 n_i.
We associate to each convex body $K$ the \emph{Hamiltonian}
\[
  H \coloneqq g_K^2.
\]
The energy surface $\{H = 1\} = \partial K$.
Where $H$ is differentiable, Hamilton's equation reads
$\dot{x}(t) = J_0 \nabla H(x(t))$.
\end{definition}

\begin{definition}[Closed characteristic]\label{def:closed-characteristic}
% Jörn: text approved (84db708)
% QC: geometric curve on bdry(K), velocity in characteristic line field.
%   Characteristic line = kernel of omega_0 restricted to T_x(bdry(K)).
%   One-dimensionality: standard fact from symplectic geometry lectures —
%   omega_0 is nondegenerate on R^4, so its restriction to a codimension-1
%   hypersurface has a 1-dimensional kernel. (Verify or flag if unfamiliar.)
%   Orientation: lambda_0 picks out a positive half-line in ell_x. The Reeb
%   vector R_x spans ell_x with lambda_0(R_x) = 1 > 0, so the positive
%   half-line is {v in ell_x : lambda_0(v) > 0}.
%   The condition lambda_0(gamma_dot) > 0 ensures: (a) gamma_dot != 0,
%   (b) gamma traverses bdry(K) in the Reeb direction (not reversed).
%   Speed is free (any positive rescaling preserves lambda_0 > 0), so any
%   reparametrization by a diffeomorphism with phi' > 0 is again a closed characteristic.
%   Ref: HKO2024 says "positive multiple of J nabla H"; our lambda_0 > 0
%   condition is equivalent (both pick the same half-line of ell_x).
% Downstream: a closed curve on bdry(K) with velocity in the positive
%   characteristic direction at each point. gamma_dot != 0 always.
%   Action A(gamma) is reparametrization-invariant (line integral of lambda_0).
Let $K$ be a convex body with smooth boundary.
At each $x \in \partial K$, the \emph{characteristic line} is
\[
  \ell_x \coloneqq \{ v \in T_x(\partial K) : \omega_0(v, w) = 0
  \;\;\forall\, w \in T_x(\partial K) \}.
\]
This is a one-dimensional subspace of $T_x(\partial K)$.
A \emph{closed characteristic} on $\partial K$ is a closed curve
$\gamma \colon [0, T] \to \partial K$ with $\gamma(0) = \gamma(T)$
whose velocity lies in the positive characteristic direction:
\[
  \dot{\gamma}(t) \in \ell_{\gamma(t)}
  \quad\text{and}\quad
  \lambda_0(\dot{\gamma}(t)) > 0
  \qquad \text{for all } t.
\]
The positivity condition orients the characteristic line and
ensures $\dot{\gamma}(t) \neq 0$.
Any reparametrization by a diffeomorphism $\varphi$ with $\varphi' > 0$
of a closed characteristic is again a closed characteristic.
The action $A(\gamma)$ is reparametrization-invariant
(it is a line integral of $\lambda_0$).
\end{definition}

\begin{definition}[Reeb orbit --- smooth]\label{def:reeb-orbit-smooth}
% Jörn: text approved (1aac085)
% QC: Reeb vector field R defined by contact normalization lambda_0(R) = 1
%   and characteristic direction omega_0(R, .) = 0 on T(bdry(K)).
%   A Reeb orbit is a periodic integral curve of R.
%   Regularity: R is a smooth vector field on a smooth manifold, so its
%   integral curves are C^infty. (Contrast polytope case: W^{1,2} only,
%   because velocity jumps at edges.)
%   Reeb orbits trace the same geometric curves as closed characteristics,
%   but carry a canonical parametrization (fixed speed).
% Downstream: R is the unique vector field on bdry(K) with
%   lambda_0(R) = 1 and omega_0(R, v) = 0 for all v in T(bdry(K)).
%   Reeb orbit gamma satisfies gamma_dot(t) = R_{gamma(t)}.
%   Polytope version: Definition~\ref{def:reeb-orbit-polytope}.
Let $K$ be a convex body with smooth boundary.
The restriction $\alpha \coloneqq \lambda_0|_{\partial K}$ is a contact form.
The \emph{Reeb vector field} $R$ on $\partial K$ is the unique vector field
satisfying
\[
  \alpha(R_x) = 1
  \quad\text{and}\quad
  \omega_0(R_x, v) = 0 \;\;\forall\, v \in T_x(\partial K),
  \qquad \text{at each } x \in \partial K.
\]
A \emph{Reeb orbit} is a periodic integral curve of $R$: a curve
$\gamma \colon [0, T] \to \partial K$ with $T > 0$,
$\gamma(0) = \gamma(T)$, and $\dot{\gamma}(t) = R_{\gamma(t)}$.

Every Reeb orbit traces the same geometric curve as a closed characteristic,
with the canonical parametrization fixed by
$\lambda_0(\dot{\gamma}) = 1$.
\end{definition}

\begin{definition}[Action]\label{def:action}
% Jörn: text approved (1aac085)
% QC: A(gamma) = integral of lambda_0 along gamma = (1/2) integral <J_0 gamma, gamma_dot>.
%   A > 0 for orbits traversed in the correct (Reeb) direction.
%   A is reparametrization-invariant (as a line integral of a 1-form).
% Downstream: A(gamma) = 0.5 * integral_0^T dot(J_0*gamma(t), gamma_dot(t)) dt.
%   Test: A(circle of radius r in q1-p1 plane) = pi * r^2.
The \emph{action} of a closed curve $\gamma \colon [0, T] \to \mathbb{R}^4$
with $\gamma(0) = \gamma(T)$ is
\[
  A(\gamma) \coloneqq \int_\gamma \lambda_0
  = \frac{1}{2} \int_0^T \langle J_0\, \gamma(t),\, \dot{\gamma}(t) \rangle \, dt.
\]
\end{definition}

\begin{remark}[Geometric meaning of action]\label{rem:action-geometric}
% Jörn: text approved (1aac085)
% Downstream: action = symplectic area of any disc bounded by gamma
%   (when gamma bounds a disc). By Stokes: integral of lambda_0 = integral of omega_0
%   over any filling disc.
By Stokes' theorem, if $\gamma = \partial D$ for a disc $D \subset \mathbb{R}^4$,
then $A(\gamma) = \int_D \omega_0$.
The action is the symplectic area enclosed by $\gamma$.
\end{remark}

\begin{lemma}[Action minima coincide]\label{lem:action-minima-coincide}
% Jörn: text approved (1aac085)
% QC: proof uses (1) every Reeb orbit is a closed characteristic,
%   (2) every closed characteristic can be reparametrized into a Reeb orbit,
%   (3) action is reparametrization-invariant (line integral of lambda_0).
%   Step (2) requires lambda_0(gamma_dot) != 0 along the characteristic.
%   Why: gamma_dot lies in the characteristic line ell_x, which is spanned
%   by R_x (the Reeb vector). Since alpha(R_x) = 1 != 0, any nonzero
%   vector in ell_x has nonzero lambda_0-evaluation. So we can rescale
%   speed to achieve lambda_0(gamma_dot) = 1, producing a Reeb orbit.
% Downstream: justifies defining c_EHZ via closed characteristics
%   and computing it via Reeb orbits interchangeably.
Let $K$ be a convex body with smooth boundary. Then
\begin{multline*}
  \min\{ A(\gamma) : \gamma \text{ is a closed characteristic on } \partial K \}
  \\= \min\{ A(\gamma) : \gamma \text{ is a Reeb orbit on } \partial K \}.
\end{multline*}
\end{lemma}

\begin{proof}
Every Reeb orbit is a closed characteristic
(Definition~\ref{def:reeb-orbit-smooth}), so the left-hand side is
${\leq}$ the right-hand side.
Conversely, any closed characteristic can be reparametrized into a Reeb
orbit (choose the parametrization with
$\lambda_0(\dot{\gamma}) = 1$).
Since $A(\gamma) = \int_\gamma \lambda_0$ is a line integral of a
$1$-form, it is reparametrization-invariant
(Definition~\ref{def:action}).
Hence the two minima coincide.
\end{proof}

\begin{definition}[EHZ capacity]\label{def:ehz-capacity}
% Jörn: text approved (1aac085)
% QC: c_EHZ(K) > 0 (positive, by thm:orbit-existence-smooth).
%   c_EHZ(K) = min (not just inf) by thm:orbit-existence-smooth.
%   Defined via closed characteristics (smooth case), extends to all convex bodies.
% Downstream: c_EHZ(K) = minimum action over all closed characteristics on bdry(K).
%   For polytopes: equivalent to minimum over generalized Reeb orbits
%   (Definition~\ref{def:reeb-orbit-polytope}), since action is
%   reparametrization-invariant and the contact normalization just picks
%   a canonical parametrization.
The \emph{Ekeland--Hofer--Zehnder capacity} of a convex body $K$ with smooth
boundary is
\[
  c_{\mathrm{EHZ}}(K) \coloneqq \min\bigl\{ A(\gamma) :
  \gamma \text{ is a closed characteristic on } \partial K \bigr\}.
\]
The minimum exists and is positive (Theorem~\ref{thm:orbit-existence-smooth}).
The definition extends to all convex bodies by approximation;
for polytopes, the minimum is attained by a generalized Reeb orbit
(Definition~\ref{def:reeb-orbit-polytope}).
\end{definition}

\begin{lemma}[Symplectic capacity axioms]\label{lem:capacity-axioms}
% Jörn: text approved (1aac085)
% QC: standard axioms from Hofer--Zehnder. Each testable.
%   No 0-in-int(K) assumption needed: these hold for all convex bodies.
% Downstream: each axiom is a proptest. Especially (ii) and (iii).
%   (i) monotonicity: test with nested polytopes.
%   (ii) conformality: test c(alpha*K) = alpha^2 * c(K) for random alpha > 0.
%   (iii) symplectic invariance: test c(phi(K)) = c(K) for random
%     symplectomorphisms phi (linear Sp(4), translations, compositions).
%   (iv) normalization: test c(ball) = pi for unit ball.
The EHZ capacity satisfies, for all convex bodies $K, K'$:
\begin{enumerate}
\item \emph{Monotonicity:} $K \subset K' \implies c_{\mathrm{EHZ}}(K) \leq c_{\mathrm{EHZ}}(K')$.
\item \emph{Conformality:} $c_{\mathrm{EHZ}}(\alpha K) = \alpha^2\, c_{\mathrm{EHZ}}(K)$
  for all $\alpha > 0$.
\item \emph{Symplectic invariance:}
  $c_{\mathrm{EHZ}}(\varphi(K)) = c_{\mathrm{EHZ}}(K)$
  for every symplectomorphism $\varphi \colon \mathbb{R}^4 \to \mathbb{R}^4$.
\item \emph{Normalization:}
  $c_{\mathrm{EHZ}}(B^4(r)) = c_{\mathrm{EHZ}}(B^2(r) \times \mathbb{R}^2) = \pi r^2$.
\end{enumerate}
\end{lemma}

% --- Capacity of symplectic products ---

\begin{proposition}[Capacity of symplectic products]%
\label{prop:capacity-symplectic-product}
% QC: uses only the four capacity axioms above plus the fact that
%   area-preserving = symplectic in 2D (omega_0 = dq wedge dp = area form).
% Downstream: correctness.tex table references this for
%   symplectic triangle products.
Let $A, B \subset \mathbb{R}^2$ be convex bodies containing
the origin in their interiors, where we view the first copy
of~$\mathbb{R}^2$ as the symplectic plane $(q_1, p_1)$ and
the second as $(q_2, p_2)$.
The \emph{symplectic product} is
\[
  A \times_S B
    \coloneqq \{(q_1, q_2, p_1, p_2) \in \mathbb{R}^4
      : (q_1, p_1) \in A,\; (q_2, p_2) \in B\}.
\]
Then
\[
  c_{\mathrm{EHZ}}(A \times_S B)
    = \min\!\bigl(\operatorname{area}(A),\;
                   \operatorname{area}(B)\bigr).
\]
\end{proposition}

\begin{proof}
We reduce to a product of disks via an area-preserving
map, then squeeze between a ball and a cylinder.

On each copy of~$\mathbb{R}^2$, the symplectic form
restricts to $dq_i \wedge dp_i$, which is the standard
area form, so area-preserving diffeomorphisms are
symplectomorphisms.

We first show $c_{\mathrm{EHZ}}(A) = \operatorname{area}(A)$
for every convex body~$A$ in~$\mathbb{R}^2$.
By the theorem of Dacorogna and
Moser~\cite{DacorognaMoser1990}, for any two bounded
domains~$\Omega_1, \Omega_2 \subset \mathbb{R}^2$ with
smooth boundaries and equal area, there exists an
area-preserving diffeomorphism
$\varphi \colon \overline{\Omega}_1 \to \overline{\Omega}_2$.
The construction solves a boundary-value problem whose
solution vanishes on~$\partial\Omega_1$, so the resulting
divergence-free vector field extends by zero outside
$\Omega_1$; its time-$1$ flow is a global area-preserving
(hence symplectic) diffeomorphism of~$\mathbb{R}^2$.
Taking $\Omega_1$ to be a smooth convex body~$A$
and $\Omega_2 = B^2(r)$ with $\pi r^2 = \operatorname{area}(A)$,
symplectic invariance gives
$c_{\mathrm{EHZ}}(A) = c_{\mathrm{EHZ}}(B^2(r)) = \pi r^2
= \operatorname{area}(A)$.
For non-smooth convex bodies, let
$A^{-\varepsilon} \subset A \subset A^{+\varepsilon}$
be smooth inner and outer parallel bodies;
monotonicity
(Lemma~\ref{lem:capacity-axioms}\,(i)) and
$\operatorname{area}(A^{\pm\varepsilon}) \to
\operatorname{area}(A)$ yield
$c_{\mathrm{EHZ}}(A) = \operatorname{area}(A)$.

Now let $r_A, r_B > 0$ satisfy
$\pi r_A^2 = \operatorname{area}(A)$ and
$\pi r_B^2 = \operatorname{area}(B)$.
The above gives global symplectomorphisms
$\varphi_1, \varphi_2 \colon \mathbb{R}^2 \to \mathbb{R}^2$
mapping $A$ to~$B^2(r_A)$ and $B$ to~$B^2(r_B)$
respectively.
Since $\omega_0 = dq_1 \wedge dp_1 + dq_2 \wedge dp_2$
decomposes over the two planes, the product
$\varphi_1 \times \varphi_2$ is a symplectomorphism
of~$\mathbb{R}^4$.
By symplectic invariance
(Lemma~\ref{lem:capacity-axioms}\,(iii)),
\[
  c_{\mathrm{EHZ}}(A \times_S B)
    = c_{\mathrm{EHZ}}\bigl(B^2(r_A) \times_S B^2(r_B)\bigr).
\]
Assume without loss of generality that $r_A \leq r_B$.
Then
\[
  B^4(r_A)
    \;\subset\; B^2(r_A) \times_S B^2(r_B)
    \;\subset\; B^2(r_A) \times \mathbb{R}^2,
\]
where the first inclusion holds because $r_A \leq r_B$
and the second because $B^2(r_B) \subset \mathbb{R}^2$.
By monotonicity and normalization
(Lemma~\ref{lem:capacity-axioms}\,(i)\&(iv)),
\[
  \pi r_A^2
    = c_{\mathrm{EHZ}}(B^4(r_A))
    \leq c_{\mathrm{EHZ}}\bigl(B^2(r_A) \times_S B^2(r_B)\bigr)
    \leq c_{\mathrm{EHZ}}(B^2(r_A) \times \mathbb{R}^2)
    = \pi r_A^2.
\]
Hence $c_{\mathrm{EHZ}}(A \times_S B)
  = \pi r_A^2
  = \min(\operatorname{area}(A), \operatorname{area}(B))$.
\end{proof}

% --- Polytope setting ---

\begin{definition}[Polytope]\label{def:polytope}
% Jörn: text approved (84db708)
% QC: K bounded, convex, 0 in int(K), intersection of half-spaces.
%   Facet count F >= 5 (at least 5 facets needed in R^4 for bounded intersection).
% Downstream: input is an h-representation {(n_i, h_i)}_{i=1}^F.
%   n_i in S^3, h_i > 0. Irredundancy means: removing any (n_i, h_i) strictly
%   enlarges the intersection.
A \emph{polytope} is a bounded convex subset
$K \subset \mathbb{R}^4$ that contains $0$ in its interior and is an
intersection of finitely many closed half-spaces.

\begin{sloppypar}
An \emph{irredundant $H$-representation} of $K$ is an unordered
finite set $\{(n_i, h_i)\}_{i=1}^{F}$ with $n_i \in S^3$,
$h_i > 0$, such that
\end{sloppypar}
\[
  K = \bigcap_{i=1}^{F} \{ x \in \mathbb{R}^4 : \langle x, n_i \rangle \leq h_i \},
\]
and removing any single inequality strictly enlarges the intersection.
\end{definition}

\begin{lemma}[Unique $H$-representation]\label{lem:h-rep-unique}
% Jörn: text approved (84db708)
% QC: the irredundant h-rep exists and is unique (as an unordered set).
%   Each (n_i, h_i) corresponds to exactly one facet.
% Downstream: polytope <-> h-rep is a bijection. Two polytopes are equal
%   iff they have the same h-rep (as unordered sets of (n_i, h_i)).
Every polytope has a unique irredundant $H$-representation,
and every irredundant $H$-representation corresponds to a polytope.
\end{lemma}

\begin{proof}
Each facet $F$ of $K$ determines a unique outward unit normal $n_F$
and height $h_F = \max_{x \in K} \langle x, n_F \rangle > 0$.
The half-space $\{x : \langle x, n_F \rangle \leq h_F\}$ is irredundant
(removing it exposes $F$).
Conversely, every irredundant inequality defines a facet.
Hence the irredundant $H$-representation consists of exactly the
facet data $\{(n_F, h_F)\}_F$, which is determined by $K$.
\end{proof}

\begin{definition}[Facets, normals, heights]\label{def:facets}
% Jörn: text approved (1aac085)
% QC: F_i is the set of points where inequality i is tight.
%   F_i is a 3-dimensional polytope (a facet = maximal proper face).
%   Irredundancy guarantees each F_i is genuinely 3-dimensional (not
%   lower-dimensional): if F_i were lower-dimensional, the inequality
%   would be redundant (implied by the other inequalities).
%   n_i is unit outward normal, h_i > 0 is the support function value.
% Downstream: facet F_i = { x in K : dot(x, n_i) = h_i }.
%   F is the number of facets (= number of inequalities in h-rep).
%   h_i = h_K(n_i). n_i perpendicular to F_i, pointing outward.
Let $K$ have irredundant $H$-representation
$\{(n_i, h_i)\}_{i=1}^{F}$.
Each facet of $K$ is $3$-dimensional (irredundancy prevents degeneracy).
The \emph{facets} are
\[
  F_i \coloneqq \{ x \in K : \langle x, n_i \rangle = h_i \},
  \qquad i = 1, \ldots, F.
\]
The vectors $n_i \in S^3$ are the \emph{outward unit normals},
the scalars $h_i = h_K(n_i) > 0$ are the \emph{heights},
and $F$ is the \emph{facet count}.
\end{definition}

\begin{definition}[Facet Reeb vectors]\label{def:reeb-vectors}
% Jörn: text approved (1aac085)
% QC: R_i = (2/h_i) J_0 n_i per correspondence.tex (our R_i = HK2017's p_i).
%   Derivation: R_i is the Reeb vector of the contact form lambda_0|_{F_i},
%   normalized so that lambda_0(R_i) = 1 on F_i.
%   On facet F_i, the Hamiltonian vector field J_0 grad(g_K^2) = R_i.
% Downstream: R_i = (2.0 / h_i) * J_0 * n_i.
%   R_i is tangent to bdry(K) (i.e., dot(R_i, n_i) = 0) because
%   dot(J_0 n_i, n_i) = omega_0(n_i, n_i) = 0.
%   |R_i| = 2/h_i (since |J_0 n_i| = |n_i| = 1).
The \emph{facet Reeb vector} on facet $F_i$ is
\[
  R_i \coloneqq \frac{2}{h_i}\, J_0\, n_i.
\]
It is the unique vector tangent to $\partial K$ along
$\operatorname{int}(F_i)$ that lies in
$\ker(\omega_0|_{n_i^\perp})$ and satisfies
$\lambda_0|_x(R_i) = 1$ for all $x \in F_i$.
\end{definition}

\begin{remark}[Reeb vector properties]\label{rem:reeb-properties}
% Jörn: text approved (1aac085)
% Downstream: all testable.
The following properties hold:
\begin{enumerate}
\item $\langle R_i, n_i \rangle = 0$ (tangent to $\partial K$).
\item $|R_i| = 2/h_i$.
\item $\omega_0(R_i, v) = 0$ for all $v \in n_i^\perp$
  (characteristic direction of $\omega_0$ restricted to the facet).
\item On $\operatorname{int}(F_i)$: $J_0 \nabla(g_K^2) = R_i$
  (Hamiltonian vector field equals Reeb vector).
\end{enumerate}
\end{remark}

\begin{definition}[Generalized Reeb orbit --- polytope]\label{def:reeb-orbit-polytope}
% Jörn: text approved (84db708)
% QC: gamma in W^{1,2}, closed, on bdry(K), velocity in conv{R_i incident}.
%   "incident" means gamma(t) in F_i.
%   conv{R_i} automatically gives lambda_0(gamma_dot) = 1 (contact normalization),
%   because lambda_0(R_i) = 1 on each facet and conv preserves this.
% Downstream: gamma: [0,T] -> bdry(K), W^{1,2}, gamma(0) = gamma(T).
%   gamma_dot(t) in convex hull of { R_i : gamma(t) in F_i } for a.e. t.
%   For simple orbits: gamma_dot = R_{sigma(k)} on each segment (pure Reeb vector).
A \emph{generalized Reeb orbit} on a polytope $K$ is a continuous curve
$\gamma \in W^{1,2}([0, T], \mathbb{R}^4)$ with $T > 0$,
$\gamma(0) = \gamma(T)$, $\operatorname{Im}(\gamma) \subset \partial K$,
such that
\[
  \dot{\gamma}(t) \in \operatorname{conv}\{ R_i : \gamma(t) \in F_i \}
  \qquad \text{for a.e.\ } t \in [0, T].
\]
\end{definition}

\begin{remark}[Contact normalization]\label{rem:contact-normalization}
% Jörn: text approved (1aac085)
% Downstream: useful sanity check in tests.
Since $\lambda_0|_x(R_i) = 1$ for every $x \in F_i$ and every
$i$, the convex combination preserves the normalization:
\[
  \lambda_0|_{\gamma(t)}(\dot{\gamma}(t)) = 1
  \qquad \text{for a.e.\ } t.
\]
In particular, the action simplifies to $A(\gamma) = T$
(the period equals the action).
\end{remark}

\begin{lemma}[Symplectic shoelace formula]\label{lem:shoelace}
% Jörn: text approved (b7a8bc0) — from \begin{lemma} to \end{proof}
% QC: direct computation, no external results needed.
%   Key steps: (a) closure kills z(0) term, (b) self-pairing vanishes
%   by omega_0 antisymmetry, (c) cross terms give the double sum.
% Downstream: used in the proof of thm:main (this section) and in
%   the simple-minimizer proof (Step~3). Corresponds to HK2017
%   Proposition~3.2.
%   Test: for a square orbit in the q_1-p_1 plane with vertices
%   (0,0), (a,0), (a,a), (0,a), the formula gives
%   2A = a^2 * omega_0(e_{p_1}, e_{q_1}) = a^2 * (-1) ... hmm,
%   need to check orientation. Better test: action = symplectic area
%   of the enclosed region.
Let $z \colon [0, T] \to \mathbb{R}^4$ be a closed curve with
piecewise-constant velocity:
$\dot{z}(t) = w_k$ on $I_k = (\tau_{k-1}, \tau_k)$ for
$k = 1, \ldots, m$, where $0 = \tau_0 < \cdots < \tau_m = T$.
Then
\begin{equation}\label{eq:shoelace}
  \int_0^T \langle -J_0\, \dot{z}(t),\, z(t) \rangle \, dt
  = \sum_{1 \leq j < i \leq m}
    |I_j|\, |I_i|\, \omega_0(w_j, w_i).
\end{equation}
\end{lemma}

\begin{proof}
Since $z$ is closed, $\sum_{k=1}^m |I_k|\, w_k = 0$.
For $t \in I_i$:
\[
  z(t) = z(0) + \sum_{j=1}^{i-1} |I_j|\, w_j
    + (t - \tau_{i-1})\, w_i.
\]
Substitute into
$\langle -J_0 w_i, z(t) \rangle$ and integrate over $[0, T]$:
\begin{align*}
\int_0^T \!\langle{-}J_0 \dot{z}, z \rangle \, dt
  &= \sum_{i=1}^m \int_{I_i}
    \!\langle -J_0 w_i,\;
      z(0) + \textstyle\sum_{j < i} |I_j| w_j
      + (t{-}\tau_{i-1}) w_i
    \rangle\, dt \\
  &= \underbrace{
      \bigl\langle -J_0 \bigl(\textstyle\sum_i |I_i| w_i\bigr),\,
      z(0) \bigr\rangle
    }_{= \, 0 \;\text{(closure)}}
  + \sum_{i=1}^m \sum_{j=1}^{i-1}
      |I_i|\, |I_j|\, \omega_0(w_j, w_i) \\
  &\quad
  + \sum_{i=1}^m
      \underbrace{\omega_0(w_i, w_i)}_{= \, 0}\;
      \cdot \tfrac{|I_i|^2}{2}.
\end{align*}
The first term vanishes because
$\sum |I_k| w_k = 0$ (closure of~$z$).
The third term vanishes because
$\omega_0$ is antisymmetric.
\end{proof}

\begin{remark}\label{rem:shoelace-name}
% Jörn: text approved (b7a8bc0)
In $\mathbb{R}^2$, $\omega_0$ reduces to the determinant and
the formula recovers the shoelace formula for the signed area of a
polygon, expressed in edge vectors rather than vertices.
\end{remark}

